%% ╔══════════════════════════════════════════════════════════════════╗
%% ║         CORIOTLAB — PLANTILLA INFORME DE ACTIVIDADES            ║
%% ║         Documento interno · flexible                            ║
%% ║         Dos variantes: PERIÓDICO y GUÍA                        ║
%% ╚══════════════════════════════════════════════════════════════════╝
%%
%% COMPILAR: xelatex -interaction=nonstopmode plantilla_actividades.tex
%%           (doble pasada — ejecutar el comando dos veces)

%% ═══════════════════════════════════════════════════════════════════
%% ZONA DE CONFIGURACIÓN — editar aquí antes de compilar
%% ═══════════════════════════════════════════════════════════════════

%% VARIANTE ACTIVA: "periodico" o "guia"
\newcommand{\TipoVariante}{periodico}

%% TIPO DE PERÍODO (solo Variante Periódico): Semanal | Quincenal | Mensual
\newcommand{\TipoPeriodo}{Semanal}

%% TÍTULO del informe (máx. 80 caracteres)
\newcommand{\DocTitulo}{[Título del informe o proyecto]}

%% NÚMERO DE DOCUMENTO — Formato: IA-AAAA-NNN
\newcommand{\DocNumero}{IA-2025-[NNN]}

%% AUTOR — nombre completo
\newcommand{\DocAutor}{[Nombre completo del autor]}

%% CARGO del autor dentro de CORIOTLAB
\newcommand{\DocCargo}{[Cargo --- CORIOTLAB]}

%% PROYECTO — nombre oficial del proyecto
\newcommand{\DocProyecto}{[Nombre del proyecto]}

%% PERÍODO REPORTADO — rango de fechas
\newcommand{\DocPeriodo}{[Período reportado: ej. "16 al 20 de junio de 2025"]}

%% FECHA DE ENTREGA — Formato: DD de mes de AAAA
\newcommand{\DocFecha}{[DD de mes de AAAA]}

%% VERSIÓN — iniciar en 1.0
\newcommand{\DocVersion}{1.0}

%% ═══════════════════════════════════════════════════════════════════

\documentclass[12pt, letterpaper]{article}

%% --- Codificación y lenguaje
\usepackage[spanish, es-nodecimaldot, es-noshorthands]{babel}

%% --- Fuentes (XeLaTeX — cargar por nombre de familia Windows)
\usepackage{fontspec}
\setmainfont{Inter}
\newfontfamily\FuenteTitulo{MuseoModerno}
\newfontfamily\FuenteCodigo{Space Mono}

%% --- Geometría — márgenes ajustados al membrete (header ~3cm, footer ~2.5cm)
\usepackage[
  top=3.5cm, bottom=3.0cm,
  left=2.5cm, right=2.5cm
]{geometry}

%% --- Color
\usepackage[table]{xcolor}
\definecolor{AzulITM}     {HTML}{102D69}
\definecolor{Magenta}     {HTML}{C14894}
\definecolor{AzulDigital} {HTML}{56ACDE}
\definecolor{GrisPizarra} {HTML}{2F2F2F}
\definecolor{GrisClaro}   {HTML}{F2F2F2}
\definecolor{GrisMedio}   {HTML}{AAAAAA}
\definecolor{GrisLinea}   {HTML}{DDDDDD}
\definecolor{Blanco}      {HTML}{FFFFFF}
\definecolor{VerdeTarea}  {HTML}{1A5C35}

%% --- Gráficos y membrete
\usepackage{graphicx}
\usepackage{tikz}
\usepackage{eso-pic}

%% --- Tablas
\usepackage{longtable}
\usepackage{booktabs}
\usepackage{array}
\usepackage{colortbl}
\usepackage{multirow}

%% --- Listas
\usepackage{enumitem}

%% --- Cajas
\usepackage[most]{tcolorbox}

%% --- Código
\usepackage{listings}

%% --- Hipervínculos
\usepackage[hidelinks]{hyperref}

%% --- Espaciado de párrafos
\usepackage{parskip}
\setlength{\parskip}{6pt}
\setlength{\parindent}{0pt}
\renewcommand{\arraystretch}{1.4}

%% -----------------------------------------------------------------
%% MEMBRETE — opacidad completa en todas las páginas
%% El membrete provee el header (logo + URL) y el footer (barra CORIOT)
%% -----------------------------------------------------------------
\AddToShipoutPictureBG{%
  \begin{tikzpicture}[remember picture, overlay]
    \node at (current page.center) {%
      \includegraphics[width=\paperwidth, height=\paperheight,
        keepaspectratio=false]%
        {../../membretes/membrete_azul.png}%
    };
  \end{tikzpicture}%
}

%% Sin encabezado ni pie de página — el membrete los provee
\pagestyle{empty}

%% -----------------------------------------------------------------
%% ESTILOS DE SECCIONES
%% -----------------------------------------------------------------
\usepackage{titlesec}
\titleformat{\section}
  {\FuenteTitulo\large\color{AzulITM}\bfseries}
  {\thesection.}{0.5em}{}
\titleformat{\subsection}
  {\FuenteTitulo\normalsize\color{GrisPizarra}\bfseries}
  {\thesubsection.}{0.5em}{}
\titlespacing*{\section}   {0pt}{1.8em}{0.7em}
\titlespacing*{\subsection}{0pt}{1.2em}{0.4em}

%% -----------------------------------------------------------------
%% MACROS DE ESTADO
%% -----------------------------------------------------------------
\newcommand{\EstadoCompletado}{%
  \textcolor{VerdeTarea}{\small\textbf{Completado}}}
\newcommand{\EstadoEnCurso}{%
  \textcolor{AzulDigital}{\small\textbf{En curso}}}
\newcommand{\EstadoPendiente}{%
  \textcolor{GrisMedio}{\small\textbf{Pendiente}}}
\newcommand{\EstadoBloqueado}{%
  \textcolor{Magenta}{\small\textbf{Bloqueado}}}

%% -----------------------------------------------------------------
%% CAJAS tcolorbox
%% -----------------------------------------------------------------
\tcbset{
  cajaNota/.style={
    colback=AzulDigital!10, colframe=AzulDigital,
    fonttitle=\FuenteTitulo\small\bfseries,
    coltitle=white, boxrule=0pt, leftrule=4pt,
    title=Nota, left=8pt, top=6pt, bottom=6pt
  },
  cajaAlerta/.style={
    colback=Magenta!8, colframe=Magenta,
    fonttitle=\FuenteTitulo\small\bfseries,
    coltitle=white, boxrule=0pt, leftrule=4pt,
    title=Advertencia, left=8pt, top=6pt, bottom=6pt
  },
  cajaResultado/.style={
    colback=AzulITM!8, colframe=AzulITM,
    fonttitle=\FuenteTitulo\small\bfseries,
    coltitle=white, boxrule=0pt, leftrule=4pt,
    title=Resultado, left=8pt, top=6pt, bottom=6pt
  },
  cajaDatos/.style={
    colback=GrisClaro, colframe=GrisLinea,
    boxrule=0.6pt, left=8pt, top=6pt, bottom=6pt
  }
}

%% -----------------------------------------------------------------
%% CÓDIGO lstlisting
%% -----------------------------------------------------------------
\lstdefinestyle{coriotcode}{
  basicstyle=\FuenteCodigo\small,
  backgroundcolor=\color{GrisClaro},
  commentstyle=\color{GrisMedio}\itshape,
  keywordstyle=\color{AzulITM}\bfseries,
  stringstyle=\color{Magenta},
  numberstyle=\tiny\color{GrisMedio},
  numbers=left, numbersep=10pt,
  breaklines=true, breakatwhitespace=false,
  frame=none, tabsize=2,
  showstringspaces=false,
  aboveskip=8pt, belowskip=8pt
}
\lstset{style=coriotcode}

%% -----------------------------------------------------------------
%% MACRO: FIRMA
%% -----------------------------------------------------------------
\newcommand{\FirmaSimple}{%
  \vspace{2.5em}
  \noindent\begin{minipage}{8cm}
    {\color{AzulITM}\rule{7cm}{0.6pt}}\\[4pt]
    {\small\bfseries\color{GrisPizarra}\DocAutor}\\[2pt]
    {\footnotesize\color{GrisMedio}\DocCargo}\\[1pt]
    {\footnotesize\color{GrisMedio}CORIOTLAB --- ITM}\\[1pt]
    {\footnotesize\color{GrisMedio}\DocFecha}
  \end{minipage}
}

%% =================================================================
%%  INICIO DEL DOCUMENTO
%% =================================================================
\begin{document}
\color{GrisPizarra}

%% -----------------------------------------------------------------
%% PORTADA — igual al Word: título directo + metadata, sin barra TikZ
%% El membrete ya provee el header con logo y la barra footer
%% -----------------------------------------------------------------

{\FuenteTitulo\LARGE\color{AzulITM}\bfseries\DocTitulo}\par\vspace{0.3em}
{\FuenteTitulo\normalsize\color{GrisPizarra}\DocProyecto}\par\vspace{1em}
{\color{GrisLinea}\rule{\linewidth}{0.5pt}}\par\vspace{0.8em}
\begin{tabular}{@{}ll@{}}
  {\small\color{GrisMedio}Autor:}   & {\small\color{GrisPizarra}\DocAutor}   \\[3pt]
  {\small\color{GrisMedio}Cargo:}   & {\small\color{GrisPizarra}\DocCargo}   \\[3pt]
  {\small\color{GrisMedio}Período:} & {\small\color{GrisPizarra}\DocPeriodo} \\[3pt]
  {\small\color{GrisMedio}Versión:} & {\small\color{GrisPizarra}\DocVersion} \\[3pt]
  {\small\color{GrisMedio}Fecha:}   & {\small\color{GrisPizarra}\DocFecha}   \\
\end{tabular}

\vspace{2em}

%% -----------------------------------------------------------------
%% === VARIANTE A: PERIÓDICO ===
%% Para reportes regulares de actividades (semanal, quincenal, mensual).
%% Estructura: Portada + Tabla de actividades + Observaciones +
%%             Próximos pasos + Firma.
%% -----------------------------------------------------------------

\section{Registro de Actividades}

%% INSTRUCCIONES DE LA TABLA:
%% - Cada fila es una actividad del período.
%% - Alternar \rowcolor{white} y \rowcolor{GrisClaro} en filas.
%% - Usar los macros de estado: \EstadoCompletado | \EstadoEnCurso |
%%   \EstadoPendiente | \EstadoBloqueado
%% - En Fecha usar formato DD/MM. En Hrs poner número entero.
%% - Para actividades bloqueadas usar "---" en Fecha y Hrs.

\setlength{\LTpre}{6pt}
\setlength{\LTpost}{6pt}

\begin{longtable}{
  >{\centering\arraybackslash}p{0.55cm}
  p{3.4cm}
  p{5.6cm}
  >{\centering\arraybackslash}p{1.6cm}
  >{\centering\arraybackslash}p{0.9cm}
  >{\centering\arraybackslash}p{2.2cm}
}
\rowcolor{AzulITM}
\color{Blanco}\FuenteTitulo\small\bfseries N\textsuperscript{o} &
\color{Blanco}\FuenteTitulo\small\bfseries Actividad &
\color{Blanco}\FuenteTitulo\small\bfseries Descripción &
\color{Blanco}\FuenteTitulo\small\bfseries Fecha &
\color{Blanco}\FuenteTitulo\small\bfseries Hrs &
\color{Blanco}\FuenteTitulo\small\bfseries Estado \\
\endfirsthead
\rowcolor{AzulITM}
\color{Blanco}\FuenteTitulo\small\bfseries N\textsuperscript{o} &
\color{Blanco}\FuenteTitulo\small\bfseries Actividad &
\color{Blanco}\FuenteTitulo\small\bfseries Descripción &
\color{Blanco}\FuenteTitulo\small\bfseries Fecha &
\color{Blanco}\FuenteTitulo\small\bfseries Hrs &
\color{Blanco}\FuenteTitulo\small\bfseries Estado \\
\endhead
%% --- FILA 1
\rowcolor{white}
1 & {[}Nombre corto de la actividad{]} &
{[}Descripción detallada de lo realizado. Máx.\ 150 caracteres. Incluir
herramientas, métodos o resultados obtenidos.{]} &
{[}DD/MM{]} & {[}H{]} & \EstadoCompletado \\
%% --- FILA 2
\rowcolor{GrisClaro}
2 & {[}Nombre corto de la actividad{]} &
{[}Descripción de actividad en progreso. Indicar porcentaje de avance
o etapa actual si es relevante.{]} &
{[}DD/MM{]} & {[}H{]} & \EstadoEnCurso \\
%% --- FILA 3
\rowcolor{white}
3 & {[}Nombre corto de la actividad{]} &
{[}Actividad planificada para el período pero aún no iniciada.
Indicar la razón del retraso si aplica.{]} &
{[}DD/MM{]} & --- & \EstadoPendiente \\
%% --- FILA 4
\rowcolor{GrisClaro}
4 & {[}Nombre corto de la actividad{]} &
{[}Actividad detenida por una dependencia externa. Describir el
impedimento y quién debe resolverlo.{]} &
--- & --- & \EstadoBloqueado \\
%% Agregar más filas según necesidad, alternando rowcolor{white} y rowcolor{GrisClaro}
\end{longtable}

\noindent{\small\color{GrisMedio}
\textbf{Total horas registradas:} {[}Total{]}\,h
\quad|\quad
\textbf{Período:} \DocPeriodo
}

\section{Observaciones}

%% Incluya: dificultades, cambios de alcance, decisiones relevantes.
%% Máximo 3 párrafos.

{\itshape\color{GrisMedio}%
Registre aquí situaciones relevantes del período: dificultades encontradas,
cambios de alcance, decisiones tomadas o cualquier información contextual
que afecte el desarrollo del proyecto. Este campo es libre pero debe ser
conciso. Máximo 3 párrafos.}

\begin{tcolorbox}[cajaNota]
  {[}Opcional: use esta caja para destacar una nota importante, alerta
  o información que requiera atención especial del lector.{]}
\end{tcolorbox}

\section{Próximos Pasos}

%% Cada ítem debe ser una acción concreta. Iniciar con verbo en infinitivo.

\begin{itemize}[label={\color{AzulITM}\textbf{>}}, leftmargin=1.8em, itemsep=4pt]
  \item {[}Acción concreta a realizar --- iniciar con verbo en infinitivo.{]}
  \item {[}Responsable y fecha límite si aplica.{]}
  \item {[}Recurso o dependencia necesaria para completar la acción.{]}
\end{itemize}

\FirmaSimple

%% -----------------------------------------------------------------
%% FIN VARIANTE A
%% -----------------------------------------------------------------


%% -----------------------------------------------------------------
%% === VARIANTE B: GUÍA TÉCNICA ===
%% Para documentar procedimientos, configuraciones o tutoriales.
%% PARA ACTIVAR: comente todo desde \section{Registro...} hasta
%%               \FirmaSimple anterior y descomente este bloque.
%% -----------------------------------------------------------------

%\section{Objetivo de la Guía}
%
%{[}Describa en 2--3 oraciones el objetivo de esta guía: qué proceso
%documenta, qué sistema configura o qué procedimiento explica.{]}
%
%\section{Requisitos Previos}
%
%\begin{itemize}[label={\color{AzulITM}\textbf{>}}, leftmargin=1.8em, itemsep=4pt]
%  \item {[}Requisito 1: hardware, software o conocimiento necesario.{]}
%  \item {[}Requisito 2: versión mínima de herramienta o librería.{]}
%  \item {[}Requisito 3: acceso o permiso requerido.{]}
%\end{itemize}
%
%\section{Pasos}
%
%\subsection{{[}Nombre del paso 1{]}}
%
%{[}Descripción del paso. Explique qué hace el usuario y por qué.{]}
%
%\begin{tcolorbox}[cajaDatos]
%\begin{lstlisting}[language=bash]
%# [Comando o código del paso 1]
%[comando-aqui]
%\end{lstlisting}
%\end{tcolorbox}
%
%\begin{tcolorbox}[cajaNota]
%  {[}Opcional: nota aclaratoria sobre el paso anterior.{]}
%\end{tcolorbox}
%
%\subsection{{[}Nombre del paso 2{]}}
%
%{[}Descripción del paso 2.{]}
%
%\begin{tcolorbox}[cajaDatos]
%\begin{lstlisting}[language=Python]
%# [Código del paso 2]
%[codigo-aqui]
%\end{lstlisting}
%\end{tcolorbox}
%
%\begin{tcolorbox}[cajaAlerta]
%  {[}Opcional: advertencia sobre un riesgo o error frecuente.{]}
%\end{tcolorbox}
%
%\section{Conclusión}
%
%{[}Resuma en 2--3 oraciones el resultado esperado al completar la guía
%y el siguiente paso natural.{]}
%
%\FirmaSimple

%% -----------------------------------------------------------------
%% FIN VARIANTE B
%% -----------------------------------------------------------------

\end{document}
